\section{隐变量与 Autoencoder}

\subsection{什么是隐变量（Latent Variable）？}

为了理解隐变量的概念，课程借用了柏拉图《理想国》中著名的\textbf{洞穴寓言}（Allegory of the Cave）：一群囚犯被束缚在洞穴中面对墙壁，他们唯一能看到的只是火光映射在墙上的物体影子。对囚犯而言，影子就是他们的现实，但真正的物体结构——他们永远无法直接观察到的部分——才是产生这些影子的根本原因。

\begin{figure}[H]
\centering
\includegraphics[width=0.85\textwidth]{slides-images/slide-10.jpg}
\caption{柏拉图的洞穴寓言类比隐变量概念\protect\footnotemark}
\end{figure}
\footnotetext{Slides 第10页。来源：柏拉图《理想国》。}

\begin{knowledgebox}{隐变量的哲学类比}
在机器学习中，我们观察到的数据就像墙上的影子——它们是某些底层因素（隐变量）的表现形式。隐变量（Latent Variables）就像洞穴中投射影子的真实物体，虽然无法直接测量，却决定了我们观测到的数据的分布。生成式建模的一个核心目标就是从观测数据中提取这些隐变量。
\end{knowledgebox}

\subsection{Autoencoder 的基本原理}

\textbf{Autoencoder}（自编码器）是一种无监督学习方法，用于从无标签的训练数据中学习低维特征表示。其核心思想是：通过一个编码器（Encoder）将高维输入数据压缩到低维隐空间（Latent Space），然后通过一个解码器（Decoder）尝试从隐空间重建原始输入。

\begin{figure}[H]
\centering
\includegraphics[width=0.85\textwidth]{slides-images/slide-13.jpg}
\caption{Autoencoder：编码器将数据 $x$ 映射到低维隐空间 $z$\protect\footnotemark}
\end{figure}
\footnotetext{Slides 第13页。}

编码器学习的映射：

$$\text{Encoder}: x \to z \quad \text{（高维} \to \text{低维）}$$

其中 $x$ 是输入数据（如一张图片），$z$ 是隐变量向量。之所以要保持 $z$ 的维度远低于 $x$，是因为我们希望网络能\textbf{压缩}数据，提炼出最本质的特征表示。

\subsection{通过重建学习隐空间}

关键问题：如何在没有标签的情况下训练编码器学习好的隐变量 $z$？

答案是：在编码器后面接一个\textbf{解码器}（Decoder），训练整个网络重建原始输入。

\begin{figure}[H]
\centering
\includegraphics[width=0.85\textwidth]{slides-images/slide-15.jpg}
\caption{完整的 Autoencoder：编码-解码-重建\protect\footnotemark}
\end{figure}
\footnotetext{Slides 第15页。}

$$\text{Decoder}: z \to \hat{x} \quad \text{（低维} \to \text{高维）}$$

训练目标是最小化输入 $x$ 与重建输出 $\hat{x}$ 之间的距离：

$$\mathcal{L}(x, \hat{x}) = \|x - \hat{x}\|^2$$

\begin{itemize}
    \item $x$：原始输入数据
    \item $\hat{x}$：经过编码-解码后的重建数据
    \item $\|\cdot\|^2$：均方误差（Mean Squared Error）
\end{itemize}

\begin{importantbox}{Autoencoder 的自监督本质}
Autoencoder 的损失函数\textbf{不需要任何外部标签}！它使用输入本身作为学习信号，通过最小化重建误差来迫使网络学习有意义的压缩表示。这使得 Autoencoder 成为一种\textbf{自监督}的无标签学习方法。
\end{importantbox}

\subsection{隐空间维度与重建质量}

隐空间的维度直接影响重建质量。Autoencoding 本质上是一种\textbf{压缩}操作：隐空间越小，信息瓶颈越大，网络被迫学习更紧凑的表示。

\begin{figure}[H]
\centering
\includegraphics[width=0.85\textwidth]{slides-images/slide-17.jpg}
\caption{隐空间维度对重建质量的影响：2D vs. 5D vs. 真实数据\protect\footnotemark}
\end{figure}
\footnotetext{Slides 第17页。}

如图所示，当隐空间为 2D 时，重建结果明显模糊；当增加到 5D 时，重建质量显著提升。这体现了表示能力与压缩程度之间的权衡。

\subsection{本章小结}

Autoencoder 通过编码器-解码器架构实现了无标签的特征学习。它将数据压缩到低维隐空间，再从隐空间重建原始输入。但传统 Autoencoder 存在一个关键问题：编码过程是\textbf{确定性的}（deterministic），给定相同输入总是产生相同的隐变量。这意味着隐空间可能不够平滑连续，无法有效用于生成新样本。这引出了下一节的 Variational Autoencoder。


